% =====================================================================
% Table 1 --- first, simplified map of the two closure axes.
% Input from section (1)introduction.tex. The assessment comments below
% were written when the table was first crafted and are kept with it.
% =====================================================================
% --- Table sketch, crafted in LaTeX (was a markdown sketch) -----------
% Assessment: the 3x2 structure is sound -- rows = financing closure,
% columns = labour-market closure -- and the F3 row matches the measured
% matrix (ALPHA/BETA-F3 = crowding out + external debt; GAMMA/DELTA-F3 =
% the Keynesian case, +2.26 % consumption / +1.78 % employment).
% Two cautions:
%  * The column labels were "fixed | flexible", which collides with the
%    paper's own vocabulary: GAMMA is the FIXED-WAGE row and ALPHA the
%    flexible-wage row, i.e. the opposite sense. The headers below name
%    the mechanism instead.
%  * The two "??" cells are filled here with the scarcity-side analogue of
%    the measured results (ALPHA-F1 ~ 0; ALPHA-F2 = full crowding out).
%    Confirm before they go in.
\begin{table}[htbp]
\centering
\small
\begin{tabular}{@{}l p{4.3cm} p{4.3cm}@{}}
\toprule
 & \multicolumn{2}{c}{\textbf{Labour-market closure}} \\
\cmidrule(lr){2-3}
 & \textbf{Factor supply fixed} (scarcity) & \textbf{Real wage fixed} (output adjusts) \\
\midrule
\textbf{Repurposing consumption} (F1)
  & pure reallocation
  & the change follows the input--output structure \\[2pt]
\textbf{Refinancing via taxes} (F2)
  & traditional crowding out
  & balanced-budget (Haavelmo) multiplier, dampened by leakage \\[2pt]
\textbf{External finance} (F3)
  & welfare loss and external debt
  & endogenous money, external debt and output expansion \\
\bottomrule
\end{tabular}
\caption{A first, simplified map of the two closure axes. The full
$5\times3$ matrix follows in Section~\ref{sec:application}.}
\label{tab:overview}
\end{table}
